% ============================================================
%  C: Como Programar -- 6ª Edição | Deitel & Deitel
%  Capítulo 13 -- O Pré-processador em C
%  EXERCÍCIOS (13.4 -- 13.10)
%  Abordagem incremental: Caps. 1 a 13
% ============================================================
\documentclass[12pt,a4paper]{article}
\usepackage[utf8]{inputenc}
\usepackage[T1]{fontenc}
\usepackage[brazil]{babel}
\usepackage{geometry}
\geometry{top=2.5cm,bottom=2.5cm,left=3cm,right=3cm}
\usepackage{parskip}\setlength{\parskip}{6pt}\setlength{\parindent}{0pt}
\usepackage{xcolor,tcolorbox,enumitem,listings,fancyhdr,titlesec,hyperref,booktabs,array,amsmath}
\tcbuselibrary{skins,breakable}

\definecolor{deitelblue}{RGB}{0,70,127}
\definecolor{deitelgray}{RGB}{240,242,245}
\definecolor{deitelgreen}{RGB}{0,110,60}
\definecolor{deitellightblue}{RGB}{220,235,250}
\definecolor{deitelred}{RGB}{160,20,20}
\definecolor{deitellorange}{RGB}{200,90,0}

\newtcolorbox{boaprática}[1][]{colback=deitellightblue,colframe=deitelblue,
  fonttitle=\bfseries\small,title={Boa prática de programação},breakable,#1}
\newtcolorbox{erroco}[1][]{colback=red!5,colframe=deitelred,
  fonttitle=\bfseries\small,title={Erro comum de programação},breakable,#1}
\newtcolorbox{respostabox}[1][]{colback=deitelgray,colframe=deitelblue!60,
  fonttitle=\bfseries\small\color{deitelblue},title={Resposta detalhada},breakable,#1}
\newtcolorbox{enunciadoBox}[1][]{colback=white,colframe=deitelblue,
  fonttitle=\bfseries,title={#1},breakable}
\newtcolorbox{saida}[1][]{colback=black!88,colframe=deitelblue,
  fontupper=\ttfamily\small\color{green!70},
  title={\small\color{white}Saída do programa},breakable,#1}

\setlist[enumerate,1]{label=\textbf{\alph*)},leftmargin=2em}
\setlist[itemize]{leftmargin=2em}

\lstset{
  basicstyle=\ttfamily\small,backgroundcolor=\color{deitelgray},
  frame=single,framerule=0.4pt,rulecolor=\color{deitelblue!40},
  breaklines=true,showstringspaces=false,
  keywordstyle=\color{deitelblue}\bfseries,
  commentstyle=\color{deitelgreen}\itshape,
  stringstyle=\color{deitelred},
  language=C,numbers=left,numberstyle=\tiny\color{gray},
  xleftmargin=18pt,xrightmargin=5pt,stepnumber=1,
}

\pagestyle{fancy}\fancyhf{}
\fancyhead[L]{\small\color{deitelblue}\textbf{C: Como Programar -- 6ª ed. | Deitel \& Deitel}}
\fancyhead[R]{\small\color{deitelblue}Cap. 13 -- Exercícios}
\fancyfoot[C]{\small\thepage}
\renewcommand{\headrulewidth}{0.4pt}

\titleformat{\section}[block]{\large\bfseries\color{deitelblue}}{\thesection.}{0.5em}{}[\titlerule]
\titleformat{\subsection}[block]{\normalsize\bfseries\color{deitelblue!80}}{}{}{}

\hypersetup{colorlinks=true,linkcolor=deitelblue}

\begin{document}

\begin{titlepage}
  \centering\vspace*{2cm}
  {\color{deitelblue}\rule{\linewidth}{2pt}}\\[0.4cm]
  {\huge\bfseries\color{deitelblue}C: Como Programar}\\[0.2cm]
  {\large\color{deitelblue}6ª Edição $\cdot$ Paul Deitel \& Harvey Deitel}\\[0.2cm]
  {\color{deitelblue}\rule{\linewidth}{2pt}}\\[1cm]
  {\LARGE\bfseries Capítulo 13}\\[0.3cm]
  {\Large\color{deitelblue!80}O Pré-processador em C}\\[0.8cm]
  \begin{tcolorbox}[colback=deitellightblue,colframe=deitelblue,width=0.85\linewidth]
    \centering{\large\bfseries Exercícios (13.4--13.10)}\\[0.2cm]
    {\normalsize Macros com argumentos e suas aplicações práticas}
  \end{tcolorbox}
\end{titlepage}

\tableofcontents\newpage

% ============================================================
\section{Exercício 13.4 -- Volume de Esfera (Macro)}
% ============================================================

\begin{respostabox}
\begin{lstlisting}
/* Ex13.4: volume_esfera.c
   Macro com um argumento + tabela de raios 1 a 10 */
#include <stdio.h>
#define PI 3.14159
#define VOLUME_ESFERA( r ) ( ( 4.0 / 3.0 ) * PI * ( r ) * ( r ) * ( r ) )

int main( void )
{
    int raio;

    printf( "%-8s%-15s\n", "Raio", "Volume" );
    printf( "----------------------\n" );

    for ( raio = 1; raio <= 10; ++raio ) {
        printf( "%-8d%-15.4f\n", raio, VOLUME_ESFERA( raio ) );
    }

    return 0;
}
\end{lstlisting}

\textbf{Saída esperada:}
\begin{saida}
Raio    Volume\\
----------------------\\
1       4.1888\\
2       33.5103\\
3       113.0972\\
4       268.0826\\
5       523.5984\\
6       904.7782\\
7       1436.7551\\
8       2144.6624\\
9       3053.6334\\
10      4188.7867
\end{saida}

\textbf{Observação importante:} note os parênteses em torno de \texttt{(r)} --
sem eles, \texttt{VOLUME\_ESFERA(2+1)} seria expandido para
\texttt{(4.0/3.0)*3.14159*2+1*2+1*2+1}, que devido à precedência da multiplicação
daria \texttt{(4.0/3.0)*3.14159*2 + 1*2 + 1*2 + 1 = 8.378 + 5 = 13.378} (errado!).
\end{respostabox}

\newpage

% ============================================================
\section{Exercício 13.5 -- Macro \texttt{SUM}}
% ============================================================

\begin{respostabox}
\begin{lstlisting}
/* Ex13.5: macro_sum.c */
#include <stdio.h>
#define SUM( x, y ) ( ( x ) + ( y ) )

int main( void )
{
    int x = 4, y = 9;
    printf( "A soma de x e y e %d\n", SUM( x, y ) );
    return 0;
}
\end{lstlisting}

\textbf{Saída:}
\begin{saida}
A soma de x e y e 13
\end{saida}

\textbf{Vantagens de macro vs.\ função:}
\begin{itemize}[noitemsep]
  \item \textbf{Sem overhead de chamada} -- a expansão ocorre em tempo de compilação.
  \item \textbf{Funciona com qualquer tipo} -- \texttt{SUM(2.5, 3.1)} também funciona.
\end{itemize}

\textbf{Desvantagens:}
\begin{itemize}[noitemsep]
  \item \textbf{Sem checagem de tipos} -- erros podem passar despercebidos.
  \item \textbf{Avaliação múltipla} -- argumentos podem ser avaliados várias vezes
  (perigoso com efeitos colaterais como \texttt{SUM(i++, j++)}).
  \item \textbf{Difícil de depurar} -- o depurador não consegue entrar em macros.
\end{itemize}
\end{respostabox}

% ============================================================
\section{Exercício 13.6 -- Macro \texttt{MINIMUM2}}
% ============================================================

\begin{respostabox}
\begin{lstlisting}
/* Ex13.6: minimum2.c */
#include <stdio.h>
#define MINIMUM2( a, b ) ( ( a ) < ( b ) ? ( a ) : ( b ) )

int main( void )
{
    double x, y;

    printf( "Digite dois valores: " );
    scanf( "%lf%lf", &x, &y );

    printf( "O menor e %.2f\n", MINIMUM2( x, y ) );
    return 0;
}
\end{lstlisting}

\textbf{Conceito:} usamos o operador ternário \texttt{? :} -- expressão que
retorna um valor, perfeita para macros (diferente de \texttt{if-else}, que é
instrução).

\begin{erroco}
  \textbf{Cuidado com efeitos colaterais:}
\begin{lstlisting}
int i = 5, j = 10;
int m = MINIMUM2( ++i, ++j );
/* expande para: ((++i) < (++j) ? (++i) : (++j))
   i sera incrementado DUAS vezes (uma na comparacao,
   outra no retorno) -- comportamento inesperado!     */
\end{lstlisting}
Em funções reais, os argumentos são avaliados uma única vez.
\end{erroco}
\end{respostabox}

\newpage

% ============================================================
\section{Exercício 13.7 -- Macro \texttt{MINIMUM3}}
% ============================================================

\begin{respostabox}
\begin{lstlisting}
/* Ex13.7: minimum3.c
   Macro definida em termos de MINIMUM2 */
#include <stdio.h>
#define MINIMUM2( a, b )    ( ( a ) < ( b ) ? ( a ) : ( b ) )
#define MINIMUM3( a, b, c ) ( MINIMUM2( MINIMUM2( a, b ), c ) )

int main( void )
{
    double x, y, z;

    printf( "Digite tres valores: " );
    scanf( "%lf%lf%lf", &x, &y, &z );

    printf( "O menor e %.2f\n", MINIMUM3( x, y, z ) );
    return 0;
}
\end{lstlisting}

\textbf{Expansão passo a passo de \texttt{MINIMUM3(5, 3, 8)}:}
\begin{enumerate}[noitemsep]
  \item \texttt{MINIMUM2(MINIMUM2(5, 3), 8)}
  \item \texttt{MINIMUM2((5 < 3 ? 5 : 3), 8)}
  \item \texttt{((5 < 3 ? 5 : 3) < 8 ? (5 < 3 ? 5 : 3) : 8)}
  \item Avaliação interna: \texttt{5 < 3} é falso $\to$ retorna 3.
  \item \texttt{(3 < 8 ? 3 : 8)} $\to$ \texttt{3}.
\end{enumerate}

\textbf{Lição:} macros podem ser \textit{composicionais} -- definir uma em termos
de outra. Esse é o princípio por trás de bibliotecas de macros como
\texttt{\#define MAX(a, b)} e \texttt{\#define MAX3(a, b, c) MAX(MAX(a, b), c)}.

\begin{boaprática}
  Note que \texttt{MINIMUM3(++i, ++j, ++k)} expandiria \texttt{++i} e \texttt{++j}
  \textit{quatro vezes} -- pelo aninhamento das macros. \textbf{Nunca passe expressões
  com efeitos colaterais para macros.}
\end{boaprática}
\end{respostabox}

\newpage

% ============================================================
\section{Exercício 13.8 -- Macro \texttt{PRINT}}
% ============================================================

\begin{respostabox}
\begin{lstlisting}
/* Ex13.8: macro_print.c */
#include <stdio.h>
#define PRINT( s ) printf( "%s\n", s )

int main( void )
{
    PRINT( "Ola, mundo!" );
    PRINT( "Macros sao uteis." );
    PRINT( "Fim do programa." );

    return 0;
}
\end{lstlisting}

\textbf{Saída:}
\begin{saida}
Ola, mundo!\\
Macros sao uteis.\\
Fim do programa.
\end{saida}

\textbf{Variação útil:} macro de log com origem usando \texttt{\_\_FILE\_\_} e
\texttt{\_\_LINE\_\_}:
\begin{lstlisting}
#define LOG( s ) printf( "[%s:%d] %s\n", __FILE__, __LINE__, s )

LOG( "Iniciando processamento" );
/* Saida: [exemplo.c:8] Iniciando processamento */
\end{lstlisting}
Esse é o padrão básico de \textit{logging frameworks} em C.
\end{respostabox}

% ============================================================
\section{Exercício 13.9 -- Macro \texttt{PRINTARRAY}}
% ============================================================

\begin{respostabox}
\begin{lstlisting}
/* Ex13.9: macro_print_array.c */
#include <stdio.h>
#define PRINTARRAY( arr, n )                       \
    do {                                            \
        int _i;                                     \
        for ( _i = 0; _i < ( n ); ++_i ) {          \
            printf( "%d ", ( arr )[ _i ] );         \
        }                                            \
        printf( "\n" );                             \
    } while ( 0 )

int main( void )
{
    int valores[] = { 10, 20, 30, 40, 50, 60 };
    int n = sizeof( valores ) / sizeof( valores[ 0 ] );

    printf( "Array: " );
    PRINTARRAY( valores, n );

    return 0;
}
\end{lstlisting}

\textbf{Truques importantes desta macro:}
\begin{itemize}[noitemsep]
  \item \textbf{\texttt{\textbackslash} no fim das linhas:} continuação da macro
  (do contrário, terminaria na primeira quebra).
  \item \textbf{\texttt{do \{ ... \} while (0)}:} idioma para encapsular múltiplas
  instruções em um bloco que pode ser usado como instrução única (com \texttt{;}).
  \item \textbf{\texttt{\_i}} em vez de \texttt{i}: convenção para evitar
  conflito com variáveis do usuário (mas macros não têm escopo próprio --
  cuidado!).
\end{itemize}

\textbf{Por que \texttt{do...while(0)}?} Para que a macro se comporte como uma
instrução simples mesmo em \texttt{if} sem chaves:
\begin{lstlisting}
if ( cond )
    PRINTARRAY( a, n );  /* expande corretamente */
else
    outraInstrucao();
\end{lstlisting}
\end{respostabox}

\newpage

% ============================================================
\section{Exercício 13.10 -- Macro \texttt{SUMARRAY}}
% ============================================================

\begin{respostabox}
\begin{lstlisting}
/* Ex13.10: macro_sum_array.c */
#include <stdio.h>
#define SUMARRAY( arr, n, total )                   \
    do {                                             \
        int _i;                                      \
        ( total ) = 0;                               \
        for ( _i = 0; _i < ( n ); ++_i ) {           \
            ( total ) += ( arr )[ _i ];              \
        }                                             \
    } while ( 0 )

int main( void )
{
    int valores[] = { 10, 20, 30, 40, 50 };
    int n = sizeof( valores ) / sizeof( valores[ 0 ] );
    int total;

    SUMARRAY( valores, n, total );

    printf( "Soma: %d\n", total );
    return 0;
}
\end{lstlisting}

\textbf{Saída:}
\begin{saida}
Soma: 150
\end{saida}

\textbf{Comparação macro vs.\ função:}

\textbf{Versão função (mais segura):}
\begin{lstlisting}
int sumArray( const int *arr, int n )
{
    int i, total = 0;
    for ( i = 0; i < n; ++i ) total += arr[ i ];
    return total;
}

/* Uso */
int total = sumArray( valores, n );
\end{lstlisting}

\textbf{Quando preferir macro a função?}
\begin{itemize}[noitemsep]
  \item \textbf{Performance crítica:} código inserido em linha, sem overhead.
  \item \textbf{Tipo genérico:} \texttt{SUMARRAY} funciona com \texttt{int},
  \texttt{double}, \texttt{long}\ldots (função precisaria de variantes).
  \item \textbf{Acesso a contexto local:} macros podem usar variáveis do
  contexto onde são expandidas.
\end{itemize}

\textbf{Quando preferir função?}
\begin{itemize}[noitemsep]
  \item \textbf{Tamanho importa:} cada uso de macro adiciona código; cada uso
  de função é só uma chamada.
  \item \textbf{Recursão:} macros não podem recursão (auto-referência é
  detectada e bloqueada).
  \item \textbf{Tipagem:} compilador valida tipos em funções; em macros, não.
  \item \textbf{Depuração:} função aparece na pilha de chamadas; macro é
  ``transparente''.
\end{itemize}

\begin{boaprática}
  Em projetos modernos, prefira \texttt{inline} (C99+) ou \texttt{static inline}
  para pequenas funções. Você ganha o melhor dos dois mundos: tipagem checada
  + expansão em linha (a critério do compilador).
\end{boaprática}
\end{respostabox}

\newpage

% ============================================================
\section*{Reflexão Final -- Capítulo 13}

\begin{tcolorbox}[colback=deitellightblue,colframe=deitelblue,
  title={\bfseries O pré-processador: ferramenta de duas faces}]

O pré-processador é uma ferramenta dupla: poderosa e perigosa.

\textbf{Onde brilha:}
\begin{itemize}[noitemsep]
  \item \textbf{Constantes simbólicas:} nomes legíveis para valores literais.
  \item \textbf{Compilação condicional:} código portátil entre Linux, Windows,
  macOS via \texttt{\#ifdef \_\_linux\_\_}, etc.
  \item \textbf{Guardas de cabeçalho:} \texttt{\#ifndef MEU\_H} previne inclusão
  múltipla.
  \item \textbf{Macros de debug:} \texttt{assert}, \texttt{LOG}, com
  \texttt{\_\_FILE\_\_} e \texttt{\_\_LINE\_\_}.
  \item \textbf{Geração de código repetitivo:} X-macros, \texttt{FOREACH}, etc.
\end{itemize}

\textbf{Onde dá problema:}
\begin{itemize}[noitemsep]
  \item \textbf{Falta de tipagem:} erros que apareceriam em compilação passam
  para tempo de execução.
  \item \textbf{Avaliação múltipla:} \texttt{MAX(f(), g())} chama \texttt{f}
  duas vezes.
  \item \textbf{Não respeita escopo:} macros são puramente textuais.
  \item \textbf{Erros confusos:} compiladores reportam erros no \textit{código
  expandido}, não no original.
\end{itemize}

\textbf{Substitutos modernos:}
\begin{itemize}[noitemsep]
  \item \texttt{const} em vez de \texttt{\#define} para constantes.
  \item \texttt{enum} para conjuntos de constantes relacionadas.
  \item \texttt{inline} (C99+) em vez de macros como pequenas funções.
  \item \texttt{\_Generic} (C11) para genéricos sem macros complexas.
  \item Templates em C++ (mas C puro ainda usa muito macro).
\end{itemize}

\textbf{Próximo capítulo:} Cap.~14 traz tópicos avançados que tecem tudo o
que aprendemos: \texttt{va\_list} (funções com número variável de argumentos),
\texttt{atexit}, manipulação de sinais, argumentos da linha de comando, e
mais. É a fronteira final entre seu conhecimento atual e a maturidade total
em C.

\end{tcolorbox}

\end{document}
